% ============================================================
%  示例：云-边-端三层系统总体架构图
%  基于 harryopo-tikz-diagram layered-arch 模板
%  Nature 无衬线风格 — 参考 paper-structure 布局经验
%
%  编译命令：xelatex -interaction=nonstopmode example-architecture-full.tex
% ============================================================

\documentclass[12pt,a4paper]{ctexart}

\usepackage{geometry}
\geometry{left=2cm,right=2cm,top=2cm,bottom=2cm}

\usepackage{fontspec}
\usepackage{tikz}
\usetikzlibrary{positioning, fit, backgrounds, arrows.meta, calc}

% ============================================================
% Nature 风格配色（与 paper-structure 一致）
% ============================================================
\definecolor{PaperBorder}{HTML}{B0C4DE}
\definecolor{PaperFill}{HTML}{FFFFFF}
\definecolor{PaperTitle}{HTML}{2C3E50}
\definecolor{PaperMuted}{HTML}{7F8C8D}
\definecolor{AccentOrange}{HTML}{D35400}
\definecolor{PaperLayerBg}{HTML}{F5F8FA}
\definecolor{PaperLayerBorder}{HTML}{D8E3EE}

% ============================================================
% 字体设置 — Nature 无衬线风格（与 paper-structure 一致）
% ============================================================
\setmainfont{Arial}
\setsansfont{Arial}
\setmonofont{Consolas}[Scale=0.90]
\setCJKsansfont{Microsoft YaHei}
\setCJKmainfont{Microsoft YaHei}

% ============================================================
%  间距常量
% ============================================================
\def\ModW{4.0cm}
\def\ModH{1.5cm}
\def\ModGap{0.55cm}
\def\LayerGap{1.6cm}
\def\LayerPad{12pt}

\begin{document}

\begin{center}
  {\Large\sffamily\bfseries\textcolor{PaperTitle}{云-边-端三层系统总体架构图}}\\[0.3em]
  {\small\sffamily\itshape\textcolor{PaperMuted}{Cloud-Edge-Device Architecture}}
\end{center}

\vspace{1em}

\begin{figure}[htbp]
\centering
\begin{tikzpicture}[
    node distance = 0pt and \ModGap,
    >=Stealth,
    line width = 0.7pt,
    every node/.style = {inner sep=0pt, outer sep=0pt},
    % ---- 模块节点：标题加粗+2行描述+技术栈 ----
    module/.style = {
        rectangle,
        draw = PaperBorder,
        fill = PaperFill,
        text = PaperTitle,
        line width = 0.8pt,
        text width = \ModW,
        minimum height = \ModH,
        font = \small\sffamily,
        align = center,
        rounded corners = 5pt,
        inner sep = 8pt,
    },
    mod-title/.style = {
        font = \small\sffamily\bfseries,
        text = PaperTitle,
    },
    mod-desc/.style = {
        font = \footnotesize\sffamily,
        text = PaperMuted,
    },
    mod-tech/.style = {
        font = \footnotesize\sffamily\ttfamily,
        text = AccentOrange,
    },
    % ---- 层容器（与 paper-structure group-box 一致） ----
    layer-box/.style = {
        draw = PaperLayerBorder,
        fill = PaperLayerBg,
        line width = 0.7pt,
        dashed,
        dash pattern = on 4pt off 3pt,
        inner sep = \LayerPad,
        rounded corners = 8pt,
    },
    layer-title/.style = {
        font = \normalsize\sffamily\bfseries,
        text = PaperTitle,
        anchor = north west,
        inner sep = 0pt,
    },
    layer-tag/.style = {
        font = \footnotesize\sffamily\itshape,
        text = PaperMuted,
        anchor = north east,
        inner sep = 0pt,
    },
    % ---- 箭头 ----
    arr-down/.style = {
        ->, >=Stealth, thick, PaperBorder,
    },
    arr-up/.style = {
        ->, >=Stealth, thick, PaperMuted,
    },
    arr-label/.style = {
        font = \footnotesize\sffamily\bfseries,
        fill = white,
        text = PaperBorder,
        inner sep = 2pt,
        outer sep = 0pt,
    },
    arr-label-rev/.style = {
        font = \footnotesize\sffamily\bfseries,
        fill = white,
        text = PaperMuted,
        inner sep = 2pt,
        outer sep = 0pt,
    },
    % ---- 图例 ----
    legend-box/.style = {
        draw = PaperLayerBorder,
        fill = PaperFill,
        line width = 0.6pt,
        rounded corners = 5pt,
        inner sep = 8pt,
        font = \footnotesize\sffamily,
        text = PaperTitle,
    },
]

    % ============================================================
    %  第一层：云层
    % ============================================================

    \node[module] (L1a) {
        {\mod-title 数据存储}\\[2pt]
        {\mod-desc 海量数据持久化}\\[1pt]
        {\mod-desc 分布式文件系统}
    };

    \node[module, right=\ModGap of L1a] (L1b) {
        {\mod-title AI 训练}\\[2pt]
        {\mod-desc 模型训练、优化}\\[1pt]
        {\mod-desc 智能算法迭代}
    };

    \node[module, right=\ModGap of L1b] (L1c) {
        {\mod-title 监控中心}\\[2pt]
        {\mod-desc 全局监控、告警}\\[1pt]
        {\mod-desc 可视化分析}
    };

    \begin{scope}[on background layer]
        \node[layer-box, fit=(L1a)(L1b)(L1c)] (box1) {};
    \end{scope}
    \node[layer-title] at ($(box1.north west)+(0,4pt)$) {云层（Cloud Platform）};
    \node[layer-tag] at ($(box1.north east)+(0,4pt)$) {数据中心 / 云服务器};

    % ============================================================
    %  第二层：边缘层
    % ============================================================

    \node[module, below=\LayerGap of L1a.south west, anchor=north west] (L2a) {
        {\mod-title 流处理引擎}\\[2pt]
        {\mod-tech Flink / DataStream}\\[1pt]
        {\mod-desc 实时数据流计算}
    };

    \node[module, right=\ModGap of L2a] (L2b) {
        {\mod-title 资源调度器}\\[2pt]
        {\mod-desc 弹性扩缩容}\\[1pt]
        {\mod-desc 任务动态分配}
    };

    \node[module, right=\ModGap of L2b] (L2c) {
        {\mod-title 缓存服务}\\[2pt]
        {\mod-tech Redis 热点缓存}\\[1pt]
        {\mod-desc 低延迟数据访问}
    };

    \begin{scope}[on background layer]
        \node[layer-box, fit=(L2a)(L2b)(L2c)] (box2) {};
    \end{scope}
    \node[layer-title] at ($(box2.north west)+(0,4pt)$) {边缘层（Edge Computing）};
    \node[layer-tag] at ($(box2.north east)+(0,4pt)$) {边缘服务器 / 网关};

    % ============================================================
    %  第三层：端层
    % ============================================================

    \node[module, below=\LayerGap of L2a.south west, anchor=north west] (L3a) {
        {\mod-title 传感器}\\[2pt]
        {\mod-desc 温度、湿度、压力}\\[1pt]
        {\mod-desc 环境数据采集}
    };

    \node[module, right=\ModGap of L3a] (L3b) {
        {\mod-title 移动设备}\\[2pt]
        {\mod-desc 智能手机、平板}\\[1pt]
        {\mod-desc 用户交互终端}
    };

    \node[module, right=\ModGap of L3b] (L3c) {
        {\mod-title IoT 网关}\\[2pt]
        {\mod-desc 协议转换、设备接入}\\[1pt]
        {\mod-desc 边缘数据汇聚}
    };

    \begin{scope}[on background layer]
        \node[layer-box, fit=(L3a)(L3b)(L3c)] (box3) {};
    \end{scope}
    \node[layer-title] at ($(box3.north west)+(0,4pt)$) {端层（Edge Devices）};
    \node[layer-tag] at ($(box3.north east)+(0,4pt)$) {终端设备 / 感知节点};

    % ============================================================
    %  层间数据流箭头
    % ============================================================

    % 端 → 边缘：数据上报
    \coordinate (e2c-start) at ($(box3.north)!0.25!(box3.north west)$);
    \coordinate (e2c-end)   at ($(box2.south)!0.25!(box2.south west)$);
    \draw[arr-up] (e2c-start) -- (e2c-end) node[midway, arr-label-rev, anchor=east, xshift=-3pt] {数据上报};

    % 边缘 → 云：聚合数据上报
    \coordinate (c2cl-start) at ($(box2.north)!0.25!(box2.north west)$);
    \coordinate (c2cl-end)   at ($(box1.south)!0.25!(box1.south west)$);
    \draw[arr-up] (c2cl-start) -- (c2cl-end) node[midway, arr-label-rev, anchor=east, xshift=-3pt] {聚合数据上报};

    % 云 → 边缘：模型/结果下发
    \coordinate (cl2c-start) at ($(box1.south)!0.25!(box1.south east)$);
    \coordinate (cl2c-end)   at ($(box2.north)!0.25!(box2.north east)$);
    \draw[arr-down] (cl2c-start) -- (cl2c-end) node[midway, arr-label, anchor=west, xshift=3pt] {模型/结果下发};

    % 边缘 → 端：控制指令
    \coordinate (c2e-start) at ($(box2.south)!0.25!(box2.south east)$);
    \coordinate (c2e-end)   at ($(box3.north)!0.25!(box3.north east)$);
    \draw[arr-down] (c2e-start) -- (c2e-end) node[midway, arr-label, anchor=west, xshift=3pt] {控制指令};

    % ============================================================
    %  图例（嵌入 tikzpicture 内部，右下角）
    % ============================================================

    \node[legend-box, anchor=north east] (legend) at ($(box3.south east)+(0,-0.8cm)$) {
        \begin{tabular}{@{}c@{\hspace{6pt}}l@{\hspace{14pt}}c@{\hspace{6pt}}l@{}}
            \textcolor{PaperBorder}{\rule[1pt]{14pt}{1.6pt}{{\scriptsize$>$}}} & {\sffamily\bfseries\small 指令/结果下发} &
            \textcolor{PaperMuted}{\rule[1pt]{14pt}{1.6pt}{{\scriptsize$>$}}} & {\sffamily\bfseries\small 数据上报}
        \end{tabular}
    };

\end{tikzpicture}
\caption{\sffamily 云-边-端三层系统总体架构图}
\label{fig:cloud-edge-device-arch}
\end{figure}

\end{document}
